\documentclass[12pt]{article}
\usepackage[utf8]{inputenc}
\usepackage[T1]{fontenc}
\usepackage{hyperref}
\title{Southeastern Ixtl\'an Zapotec Language Resources for K-12 Curriculum Development}
\author{Compiled by Claude Sonnet 5.5 (claude-sonnet-5-5) for Andr\'as Kornai}
\date{}
\begin{document}
\maketitle

\section{Introduction}
Southeastern Ixtl\'an Zapotec (ISO 639-3 zpd, Glottocode sout3005) is an Otomanguean language of Oaxaca, Mexico. Other names listed in the sources are Latuvi Zapotec, Yaves\'ia Zapotec and Zapoteco del Sureste de Ixtl\'an. Wikitongues gives it an EGIDS status of 6b (threatened) and says it is unwritten. Very little material exists. This report found no school materials, grammar or dictionary in the search results.

\section{Primary Educational Resources}
No primers, textbooks or school programmes for this language turned up in the searches.

\section{Linguistic and Reference Materials}
\begin{itemize}
\item Glottolog entry (classification and bibliography pointers): \url{https://glottolog.org/resource/languoid/id/sout3005}
\item AILLA (University of Texas) catalog entry, Zapotec, Southeastern Ixtl\'an: \url{https://islandora-ailla.lib.utexas.edu/islandora/object/ailla%3A247555}
\item Wikitongues language page: \url{https://wikitongues.org/languages/zpd/}
\end{itemize}

\section{Religious and Community Resources}
\begin{itemize}
\item Joshua Project people-group profile (Zapoteco, Ixtlan Sureste), which lists Scripture translation as started. Its narrative text looks generic, so verify the details elsewhere: \url{https://joshuaproject.net/people_groups/print/15012/MX}
\end{itemize}

\section{Digital and Audio Resources}
\begin{itemize}
\item Global Recordings Network language page: \url{https://globalrecordings.net/en/language/zpd}
\item Global Recordings Network page for the Ixtepeji variety, with evangelistic and Bible-teaching audio. The two GRN pages disagree about whether recordings exist, so check both: \url{https://glorecnet-website.k.globalrecordings.net/en/language/3402}
\end{itemize}

\section{Curriculum Development Context}
Materials in the language appear to be almost absent. A curriculum would first need community-led orthography work, documentation, and the creation of basic readers. The AILLA archive is the likeliest source of existing recordings.
\end{document}